% 47-19(47쪽) — 곡선 y=x^3 과 점 P 에서의 접선, 그 접선에 수직이고 P 를 지나 y축과 점 Q 에서 만나는 직선. 스캔 화질이 낮아 TikZ 로 다시 그림.
% 원문에 있는 것만: 라벨 O, x, y, P, Q, y=x^3 과 P 의 직각 표시. 원문에 눈금 숫자가 없어 넣지 않았다(P 의 좌표도 적지 않는다).
% x·y 단위를 같게 두어 두 직선의 직각이 실제로 직각으로 보이게 했다.
\begin{tikzpicture}[x=0.40cm, y=0.40cm, line width=0.5pt]
  \draw[->, >=stealth, line width=0.7pt] (-1.9,0) -- (3.0,0) node[below=1pt] {$x$};
  \draw[->, >=stealth, line width=0.7pt] (0,-3.0) -- (0,3.6) node[left=1pt] {$y$};
  \node[above left=-2pt and -1pt, font=\small] at (0,0) {$\pt{O}$};
  % 접선(기울기 3)과 접선에 수직인 직선(기울기 -1/3) — 접점 P(1,1)
  \draw[line width=0.5pt] (-0.3,-2.9) -- (1.45,2.35);
  \draw[line width=0.5pt] (-1.7,1.9) -- (2.8,0.4);
  % 곡선 y=x^3
  \draw[line width=0.6pt, smooth] plot coordinates {
    (-1.40,-2.744) (-1.30,-2.197) (-1.20,-1.728) (-1.10,-1.331) (-1.00,-1.000)
    (-0.85,-0.614) (-0.70,-0.343) (-0.55,-0.166) (-0.40,-0.064) (-0.20,-0.008)
    (0.00,0.000) (0.20,0.008) (0.40,0.064) (0.55,0.166) (0.70,0.343) (0.85,0.614)
    (1.00,1.000) (1.10,1.331) (1.20,1.728) (1.28,2.097) (1.35,2.460)};
  % P 의 직각 표시(접선 방향 (1,3)/√10 · 수직선 방향 (-3,1)/√10)
  \draw[line width=0.4pt] (1.0791,1.2372) -- (0.8419,1.3163) -- (0.7628,1.0791);
  % 라벨(원문 자리 그대로)
  \node[right=2pt, font=\small] at (1.05,0.95) {$\pt{P}$};
  % y 단위가 작아 교점(y=4/3)과 원점 사이에 라벨 둘이 안 들어간다. Q 는 교점 바로 위 오른쪽(선·곡선·접선이 없는 쪽)에 둔다.
  \node[anchor=south west, inner sep=1.5pt, font=\small] at (0.10,1.3333) {$\pt{Q}$};
  \node[right=1pt, font=\small] at (0.3,3.15) {$y=x^3$};
\end{tikzpicture}
